\chapter{An uncertain customer identity changes the economic question}
\label{gap29:observation}
A customer table and an outcome table do not become one economic sample merely
because a matching algorithm returns a name. A wrong match assigns somebody
else's outcome to the intervention. An unmatched record removes an outcome
from the comparison. These are different errors, and either can change the
case for spending.

The study generates observations for 1,000 independent people. Intervention
assignment follows $A_i\sim\operatorname{Bernoulli}(1/2)$, independently of
the characteristic $X_i\sim U[-1,1]$ and disturbance
$\epsilon_i\sim U[-1,1]$. Outcomes satisfy
\[
 Y_i=100+30A_i+10X_i+20\epsilon_i.
\]
The population effect is 30 illustrative units and all outcomes lie in
$[70,160]$. Names and birth-day categories create deliberately artificial
collisions. Independently assigned row numbers and shuffled records carry
no cross-table identity. The analyst sees intervention and characteristics
in one table, and identifiers and outcome in the other. True links remain
in a separate evaluator table.

We corrupt exactly zero, 10\% or 30\% of outcome records. Corruption alternates
between a missing name, a one-character edit, another person's identifiers
and a conflicting duplicate outcome. Records are chosen either independently,
in proportion to $1+0.8X$, or with greater probability for larger outcomes.
There are 200 independent replications of each of the nine cases.
This construction tests consequences of specified errors; it does not
estimate the prevalence or form of errors in a bank.

\section{What matching can and cannot establish}
The simplest analysis accepts a unique exact name-and-date match. A fuzzy
alternative permits a one-character name difference at the same date and
uses date candidates when the name is missing. Both abstain when a candidate
is ambiguous or the same outcome is proposed for several people. Neither
uses the true identity or the intervention outcome to choose a match.

With outcome-dependent corruption at 30\%, exact matching retains 65.6\%
of people and assigns a wrong outcome to 1.56\% of the original population.
Its effect estimate has bias $-0.560$ (Monte Carlo standard error $0.071$);
its normal interval covers the population effect in 93.5\% of replications.
The fuzzy rule retains 62.8\%, with wrong matches for 0.477\% of the initial
population. Its bias is $-0.071$ (Monte Carlo standard error $0.079$) and
coverage is 94.5\%. These fractions describe this corruption process.
They do not certify a general matching rule or a nominal coverage guarantee
under outcome-dependent selection.

A candidate list has a separate interpretation. If it contains the true
outcome, its smallest and largest outcomes bound that person's observed
contribution. When no candidate exists, the declared support supplies the
bounds. A narrow candidate list is not proof that it includes the truth.
In particular, copied identifiers can exclude the true match while producing
an apparently convincing candidate.

\section{A bound whose assumption is visible}
Let $n_a$ be the initial number assigned to arm $a$, and define
$w_i=1/n_1$ for treated people and $w_i=-1/n_0$ for controls.
For each person let $[l_i,u_i]$ bound the weighted contribution using candidate
outcomes. Let $[L_i,U_i]$ be the wider weighted range obtained from $[70,160]$.
Suppose at most $K$ people have a candidate list that excludes their true
outcome. Then the full-data realised difference in means $D$ satisfies
\[
 \sum_i l_i-\sum_{j=1}^{K} d_{(j)}
 \ \le D\le\
 \sum_i u_i+\sum_{j=1}^{K} e_{(j)},\qquad
 d_i=l_i-L_i,\quad e_i=U_i-u_i,
\]
where the parenthesised indices order the nonnegative corrections from
largest to smallest. To see the lower bound, each correct candidate list
contributes at least $l_i$. Replacing an incorrect list can lower that
contribution by at most $d_i$. At most $K$ replacements cannot reduce the sum
more than its $K$ largest possible corrections. The upper argument is
identical. Relaxing one-to-one assignment constraints makes the result
conservative; no claim of sharpness is needed.

Here the analyst is supplied the corruption budget $K$, not the locations
or types of the errors. That budget bounds the number of excluded true
candidates in the specified construction. In a bank it would require an
external audit or be varied explicitly as a sensitivity assumption.
Matching confidence does not supply that information.

Sampling uncertainty is additional to linkage uncertainty. For independent
outcomes with range 90, the centred log moment-generating function has second
derivative at most $90^2/4$, since it is a variance under an exponentially
tilted distribution on the same support. Integrating twice gives the bound
$\log E\exp\{\lambda(Y-EY)\}\le\lambda^2 90^2/8$. Exponential Markov
and minimisation over $\lambda$ give a one-tail probability
$\exp(-2n_a t^2/90^2)$ for a group-mean error of size $t$.
Applying this to the four tails of the two group means yields a simultaneous
95\% envelope of radius
\[
 h=90\left\{\sqrt{\frac{\log 80}{2n_1}}+
             \sqrt{\frac{\log 80}{2n_0}}\right\}
\]
for the difference. Enlarging the linkage interval by $h$ therefore
covers the population contrast under the stated random assignment,
independent sampling, support and exclusion-budget assumptions.
The realised arm counts are retained; the argument is conditional on them.

\section{The price of honest ignorance}
\begin{figure}[htbp]
\centering
\includegraphics[width=.95\textwidth]{scientific-gap-closure/2026-10-29/linkage-width.pdf}
\caption{Outcome-dependent record corruption. Matching intervals remain
narrow while the explicitly justified exclusion-budget envelope widens.
The conditional candidate calculation relies on candidate inclusion;
the broader envelope permits a supplied number of exclusions.}
\end{figure}
The corruption-aware set contains the full-data realised contrast in every
one of the 1,800 simulations. An independent exhaustive six-person example
also verifies the largest-correction arithmetic. Neither check supplies the
corruption budget in an actual application; the derivation shows where that
premise enters.

At 30\% outcome-dependent corruption, the average population interval has
width 119.1, compared with about 4.1 for the fuzzy matched analysis. It covers
the population effect in all 200 replications. A Wilson interval for that
observed coverage fraction has a lower endpoint about 98.1\%; zero observed
failures is not a universal empirical guarantee. The narrower conditional
candidate envelope also covers all 200 times in this example, but lacks the
exclusion-budget argument when true candidates are absent.

The economically useful repair is to narrow justified uncertainty: improve
record identity, validate a tighter outcome range, or audit enough links to
defend a smaller exclusion budget. Forcing every customer to one outcome
merely conceals the information problem. An analysis that cannot support a
spending decision can still show which information is worth buying.
